\section{Evidencia de Scaling y Generalization de RT-1}

La success rate promedio en las tareas de entrenamiento no indica si el robot puede entrar a una cocina nueva, ignorar un distractor o combinar instrucciones no vistas. Por eso, RT-1 se evalúa a lo largo de cuatro ejes: seen tasks, unseen tasks, distractors y backgrounds, y además se prueba su desempeño cuando varios factores de variación aparecen al mismo tiempo.

\subsection{Visual y Semantic Diversity}

La sección anterior solo demostró que RT-1 puede ejecutar control multitarea bajo restricciones de tiempo real; aún no respondía si estas capacidades dependen de la cocina de entrenamiento, de objetos fijos y de instrucciones con plantilla. Esta sección primero descompone la generalization en cuatro cambios de distribución distinguibles; al leer la figura, conviene comparar por separado las condiciones seen, unseen, distractor y background, en lugar de fijarse solo en la barra más alta o en el promedio global. Estos resultados respaldan la ventaja frente a los baseline, pero no demuestran que el modelo haya adquirido comprensión semántica de mundo abierto.

Del lado izquierdo de la slide 17, cuatro grupos de barras comparan RT-1, GATO, BC-Z y BC-Z XL; del lado derecho se muestran variaciones de fondo, de objetos y de distractores. RT-1 es en conjunto más fuerte en las cuatro condiciones, pero las condiciones unseen / background siguen claramente por debajo de las seen tasks, lo que indica que ser el «mejor baseline» no equivale a haber resuelto la generalización en mundo abierto.

\lecturefigure{slide-17-rt1-visual-semantic-diversity.jpg}{Tasa de éxito de RT-1 ante variaciones seen, unseen, distractor y background}{Diapositiva V017 recuperada de la grabación oficial; Brohan et al., 2022}

\begin{knowledgebox}{Lectura de la figura: cada barra mide un modo de falla distinto}
Las seen tasks evalúan principalmente el ajuste a la distribución de entrenamiento; las unseen tasks evalúan la composición de lenguaje y conducta; los distractors evalúan si objetos irrelevantes desvían la atención; los backgrounds evalúan los cambios en la distribución visual. Los cuatro rubros no pueden reducirse a un solo promedio, porque cada architecture component contribuye de manera distinta a cada eje.
\end{knowledgebox}

\subsection{La variación de múltiples factores es más difícil que la de uno solo}

La figura anterior aísla los factores de variación uno por uno, lo que facilita ubicar en qué tipo de shift se degrada el modelo; sin embargo, en un despliegue real rara vez cambia una sola variable a la vez. Por ello, esta sección pasa al compound shift: la pregunta es si la ventaja ante un solo factor se mantiene cuando cambian simultáneamente el fondo, la disposición, la posición de los objetos, el punto de vista y el lenguaje. Al leer la figura, primero compara el descenso progresivo de level 1 a level 3 y después observa las diferencias concretas en la unseen kitchen; esta demostración indica que la variación conjunta es más difícil, pero unos cuantos casos exitosos no deben tomarse como evidencia de que el modelo se transfiere a cualquier entorno nuevo.

La slide 18 coloca lado a lado los resultados en barras y el video de la unseen kitchen: el robot no solo cambia de fondo, sino que puede enfrentar al mismo tiempo nuevas posiciones de objetos, otra superficie de mesa, variaciones de lenguaje y diferencias de punto de vista. En las stacked bars, los level 1/2/3 indican cada vez más factores que cambian simultáneamente; aunque RT-1 va a la cabeza, su tasa de éxito en la dificultad más alta sigue siendo limitada.

\lecturefigure{slide-18-rt1-variation-unseen-kitchen.jpg}{Generalization multifactorial: cambio simultáneo de fondo, disposición y tarea en una unseen kitchen}{Diapositiva V018 recuperada de la grabación oficial; Stanford CS25 Lecture 15}

El segundo demo usa variantes en lenguaje natural; por ejemplo, ya no repite la plantilla de entrenamiento, sino que pide al robot que actúe mediante descripciones más libres. Esto muestra el valor del language conditioning, pero también deja ver que la semántica de la tarea sigue dependiendo de la skill coverage existente: el modelo puede entender otra forma de decir lo mismo, pero no puede crear con el lenguaje capacidades físicas que los datos de entrenamiento nunca incluyeron.

\lecturefigure{slide-19-rt1-language-variation-demo.jpg}{Demostración de variación de lenguaje: ejecución de tareas con distinta redacción bajo variación visual multifactorial}{Diapositiva V019 recuperada de la grabación oficial; Stanford CS25 Lecture 15}

\teachervoice{El ponente presenta estos videos como complemento de las gráficas cuantitativas, no como una «selección de los mejores casos». Insiste en que la generalization multifactorial es mucho más difícil que las pruebas de una sola variable; que un robot funcione algunas veces en una cocina nueva no significa que ya haya aprendido cualquier entorno doméstico.}

\subsection{Mezcla de datos entre Distribution}

RT-1 también incorpora datos de simulation / sim-to-real, de bin-picking y de distintas embodiment. En la figura, a la izquierda están las fuentes de datos y a la derecha las barras reportan la mejora relativa que aportan a las targeted evaluations. La conclusión es que un modelo de alta capacidad puede extraer información compartida de datos heterogéneos, no que todos los datos produzcan transferencia positiva incondicional.

\lecturefigure{slide-20-rt1-diverse-data-distributions.jpg}{Entrenamiento con múltiples data distributions: everyday robots, simulation y bin picking}{Diapositiva V020 recuperada de la grabación oficial; Brohan et al., 2022}

\begin{warningbox}{La Data interoperability tiene límites}
Cada robot tiene cámaras, espacios de acción y dinámicas distintos. Sin un embodiment identifier, action normalization o una semántica compartida, concatenar los datos sin más genera conflictos. El resultado de RT-1 es que «parte de los datos heterogéneos puede ayudar a distribuciones específicas», no que «cualquier robot dataset puede fusionarse directamente».
\end{warningbox}

\subsection{La Task Diversity importa más que el Raw Size}

En la data scaling slide, el eje horizontal es la proporción de datos conservada y el eje vertical es la success rate; los puntos azules son las seen tasks y los negros, la generalization. La flecha verde indica reducir la cantidad total de datos; la flecha morada indica reducir la task diversity conservando una mayor cantidad de episodes. Esta última provoca una caída más pronunciada, lo que muestra que repetir la misma tarea no sustituye cubrir más tareas.

\lecturefigure{slide-21-rt1-data-scaling.jpg}{Scaling with data: reducir la task diversity perjudica más la generalización que reducir la misma cantidad de episodes}{Diapositiva V021 recuperada de la grabación oficial; Brohan et al., 2022}

\begin{importantbox}{Las cuatro «escalas» de los datos robóticos}
\begin{enumerate}
  \item \textbf{Episode count}: número total de trayectorias;
  \item \textbf{Task diversity}: combinaciones de habilidades y objetivos distintos;
  \item \textbf{Environment diversity}: fondo, disposición, objetos y lighting;
  \item \textbf{Embodiment diversity}: distintos robots y action space.
\end{enumerate}
«Más datos» solo tiene sentido en ingeniería si se especifica de qué escala se trata.
\end{importantbox}

\subsection{Ablation: no es un efecto exclusivo del Transformer}

La gráfica de ablation toma el RT-1 completo como punto cero; hacia abajo se representa la pérdida de tasa de éxito. GATO, BC-Z, BC-Z XL, quitar el modelo grande, quitar el preentrenamiento, cambiar a action continua/autorregresiva, quitar el history o el Transformer producen caídas distintas en los cuatro ejes de evaluación. El resultado indica que la architecture, el preentrenamiento, el history y la action representation intervienen en el desempeño, y que este no depende solo del número de parámetros.

\lecturefigure{slide-22-rt1-design-ablations.jpg}{RT-1 design ablations: impacto relativo de cada componente en los cuatro tipos de generalization}{Diapositiva V022 recuperada de la grabación oficial; Brohan et al., 2022}

\begin{knowledgebox}{Lectura de la figura: no busques solo la barra más larga}
Cada color corresponde a un evaluation axis distinto; un componente puede ayudar sobre todo en las unseen tasks y tener poco efecto en las seen tasks. La pérdida de `w/o transformer` o `w/o pretraining` no puede interpretarse por sí sola como una causal law, porque la estructura sustituta, la estabilidad del entrenamiento y la capacidad pueden cambiar al mismo tiempo. El valor de la Ablation está en ubicar los ejes frágiles, no en asignar un orden absoluto a los componentes.
\end{knowledgebox}

\subsection{Resumen de la sección}

La evidencia de RT-1 muestra que una Transformer-style token interface puede mejorar, dentro de un presupuesto de tiempo real, el control multitarea y varios ejes de generalización, aunque la capacidad sigue disminuyendo con la complejidad del escenario. La conclusión más sólida sobre los datos es que la task diversity pesa más que la simple episode repetition; la conclusión más sólida sobre el sistema es que varios componentes actúan en conjunto, y el éxito no puede atribuirse a la sola frase «usó un Transformer».
